\documentclass[10pt]{article}

\usepackage[utf8]{inputenc}

\usepackage[T1]{fontenc}

\usepackage{amsmath}

\usepackage{amsfonts}

\usepackage{amssymb}

\usepackage[version=4]{mhchem}

\usepackage{stmaryrd}

\usepackage{bbold}



\begin{document}

кция $S$ задает лежандрово многообразие в контактном пространстве $V \times V^{*} \times \mathbb{R}$ 1-струй функций на $V^{*}$ по формуле



\[

p=\partial s / \partial q, u=\langle p, q\rangle-S(q), p \in V^{*}, q \in V, u \in \mathbb{R} .

\]



В случае, когда функция $S$ выпукла, фронт этого лежандрового многообразия в $V^{*} \times \mathbb{R}$ снова является графиком выпуклой функции $u=S^{*}(p)-$ Л. П. функции $S$.



Механич. система классич. механики может быть эквивалентным образом описана либо функцией Лагранжа на пространстве касательного расслоения конфигурационного многообразия системы, либо функцией Гамильтона на пространстве кокасательного расслоения. Переход от функции Лагранжа к функции Гамильтона и обратно осуществляется при помощи послойного Л. П. ЛЕЖАНДРА ПРИСОЕДИНЕННЫЕ ФУНКЦИИ см. Сферические функции.



ЛЕЖАНДРА УРАВНЕНИЕ - см. Лежандра многочлены. ЛЕЖАНДРА ФУНКЦИОНАЛЬНЫЕ ПРЕОБРАЗОВАНИЯ - составная часть функциональной техники, общей для квантовой теории поля (евклидовой и псевдоевклидовой), квантовой статистики и теории классического случайного поля. Л. Ф. П. являются непосредственным обобшением обычных Лежандра преобразований функций, широко применяемых в механике и термодинамике: функции заменяются функционалами, обычные производные - вариационными.



Хорошо изучены преобразования Лежандра производящего функционала связных функций Грина



\[

W(a)=\ln \int D \varphi \exp S(\varphi, a),

\]



где $\int D \varphi . .-$ символ функционального интегрирования, $S-$ функционал действия общего вида



\[

S(\varphi, a)=\sum_{n} \int \ldots \int d x_{1} \ldots d x_{n} a_{n}\left(x_{1}, \ldots, x_{n}\right) \varphi\left(x_{1}\right) \ldots \varphi\left(x_{n}\right)

\]



(терминология евклидовой теории, в псевдоевклидовой $S$ заменяется на $i S$ ). Функционал $S$ задается набором симметричных по перестановкам $x$ потенциалов $a_{n}\left(x_{1}, \ldots, x_{n}\right)$, $n=1,2, \ldots$, величина $W-$ функшионал от всех потеншиалов $a$, его кратные вариационные производные по первому потенциалу $a_{1}(x)$ имеют смысл связных функций Грина в теории с действием $S$. Преобразование Лежандра $W$ по первому потенциалу $a_{1}$ называется первым (оно же - эф фективное действие), по совокупности $a_{1}$ и $a_{2}-$ вторым, по совокупности $a^{\prime} \equiv\left(a_{1}, \ldots, a_{m}\right)-$ преобразованием порядка $m$. При этом $a=\left(a^{\prime}, a^{\prime \prime}\right)$, где $a^{\prime \prime}-$ набор всех старших потенциалов $a_{n} \mathrm{c} n>m$, считающихся фиксированными параметрами. Преобразование порядка $m$ есть функционал



\[

\Gamma_{m}\left(\alpha, a^{\prime \prime}\right)=W(a)-\sum_{n=1}^{m} \alpha_{n} a_{n}

\]



c



\[

\alpha_{n} a_{n} \equiv \int \ldots \int d x_{1} \ldots d x_{n} \alpha_{n}\left(x_{1}, \ldots, x_{n}\right) a_{n}\left(x_{1}, \ldots, x_{n}\right),

\]



в к-ром вместо $a^{\prime}$ новыми и независимыми переменными считаются сопряженные с $a_{n}, n \leqslant m$, полные функции Грина



\[

\alpha_{n}\left(x_{1}, \ldots, x_{n}\right)=\delta W(a) / \delta a_{n}\left(x_{1}, \ldots, x_{n}\right)=\left\langle\varphi\left(x_{1}\right), \ldots, \varphi\left(x_{n}\right)\right\rangle,

\]



где $\langle. .$.$\rangle обозначает функциональное усреднение с весом$ $\exp S$. Из определения $\Gamma_{m}$ вытекают равенства:



\[

\delta \Gamma_{m}\left(\alpha, a^{\prime \prime}\right) / \delta \alpha_{n}\left(x_{1}, \ldots, x_{n}\right)=-a_{n}\left(x_{1}, \ldots, x_{n}\right),

\]



эквивалентные условию стационарности по $\alpha$ функционала



\[

\Phi_{m}(\alpha, a)=\Gamma_{m}\left(\alpha, a^{\prime \prime}\right)+\sum_{n=1}^{m} \alpha_{n} a_{n}

\]



при фиксированных $a$. В физич. задачах функции Грина $\alpha$ ишутся по заданному действию, то есть по $a$. На языке $\Phi_{m}$ задача определения первых $m$ функций $\alpha$ по $a$ становится вариационной, так что это - естественный язык для описа- ния аномальных решений со спонтанным нарушением симметрии (теория фазовых переходов и т. п.): $\Phi_{m}$ всегда сохраняет симметрию действия, но может при этом иметь неинвариантные точки стационарности, инвариантна лишь их совокупность в целом. Важно, что сохранение симметрии у $\Phi_{m}$ позволяет вычислять этот функционал, в отличие от самих функций Грина $\alpha$, по обычной теории возмущений. Функционалы $W$ и $\Gamma_{m}$ имеют известные диаграммные представления, причем основная масса диаграмм $\Gamma_{m}$ (для $\Gamma_{1}$ и $\Gamma_{2}-$ все) обладает свойством реберной $m$-неприводимости, то есть сохраняет связность при разрыве любых $m$ линий.



Более сложными и менее изученными объектами являются преобразования Лежандра функционала $S$-матрицы, диаграммы к-рых обладают свойствами вершинной неприводимости.



Jum.: [1] Dominicis C. de, Martin P. C., «J. Math. Phys.», 1964, v. 5, p. 14-59; [2] В а с и л в в в А. Н., Функциональные методы в квантовой теории поля и статистике, Л., 1976. А. Н. Васильев. ЛЕХАНДРОВ КОБОРДИЗМ - отношение эквивалентности между лежандровыми подмногообразиями лежандрово расслоенных контактных многообразий или, эквивалентно, между волновыми фронтами; фронт на краю многообразия лежандрово кобордантен нулю, если он высекается фронтом, лежащим во всем многообразии, и трансверсальным краю. Вычисления групп Л. к. параллельно (и несколько проще) вычислений для лагранжева кобордизма.



Лит.: [1] Арн ольд В. И., Ги венталь А. Б., вкн.: Итоги науки и техники. Современные проблемы математики. Фундаментальные направления, т. 4 , М., 1985 , с. 8-139. ЛЕЖАНДРОВ МОДУЛЬ элли птического интегра л а - см. Эллиптический интеграл.



ЛЕЖАНДРОВ ХАРАКТЕРИСТИЧЕСКИЙ КЛАСС класс в когомологиях лежандровых многообразий, определяющий инвариант лежандрова кобордизма. Известные Л. х. к.- Маслова класс и другие классы, индуцированные из когомологий грассманианов лагранжевых плоскостей (классы Маслова-Арнольда и Бореля-Фукса), а также классы, соответствующие когомологиям универсальных комплексов лежандровых особенностей.



Лит.: [1] Ар н о л ьд В. И., «Успехи математич. наук», 1983, т. 38, в. 2, с. 77-147; [2] Vassilye v V. A., Lagrange and Legendre charac$\begin{array}{ll}\text { teristic classes. N. Y., } 1988 . & \text { B. } \text { A. Bacuлbee. }\end{array}$



ЛЕЖАНДРОВА ОСОБЕННОСТЬ; Н о р м а ь Н а Я фор м а - наиболее простой вид, к которому можно привести росток лежандрового отображения лежандровой эквивалентностью. Пусть $A_{\mu}, D_{\mu}, E_{6}, E_{7}, E_{8}-$ ростки лежандровых отображений, заданных следуюшими производящими семействами:



\[

\begin{gathered}

A_{\mu}, \mu \geqslant 1: x^{\mu+1}+q_{\mu} x^{\mu-1}+\ldots+q_{2} x+q_{1}=0, \\

D_{\mu}, \mu \geqslant 4: x^{2} y+y^{\mu-1}+q_{\mu} x+q_{\mu-1} y^{\mu-2}+\ldots+q_{2} y+q_{1}=0, \\

E_{6}: x^{3}+y^{4}+q_{6} x y^{2}+q_{5} x y+q_{4} y^{2}+q_{3} x+q_{2} y+q_{1}=0, \\

E_{7}: x^{3}+x y^{3}+q_{7} x^{2} y+q_{6} x^{2}+q_{5} x y+q_{4} y^{2}+q_{3} x+q_{2} y+q_{1}=0, \\

E_{8}: x^{3}+y^{5}+q_{8} x y^{3}+q_{7} x y^{2}+q_{6} y^{3}+q_{5} x y+q_{4} y^{2}+q_{3} x+q_{2} y+q_{1}=0 .

\end{gathered}

\]



Лежандрово отображение общего положения $(l-1)$-мерного многообразия в $l$-мерное при $l \leqslant 6$ устойчиво в окрестности любой точки и эквивалентно отображению $A_{\mu}, D_{\mu}$, $E_{\mu}, \mu \leqslant l$. При $l \geqslant 7$ лежандровы отображения могут быть неустойчивы и иметь (при $l \geqslant 10-$ функциональные) модули. $\quad$ М. Э. Казарян. ЛЕЖАНДРОВА ЭКВИВАЛЕНТНОСТЬ - см. Лежандрово отображение.



ЛЕЖАНДРОВО МНОГООБРАЗИЕ - инТегральНое мНогообразие максимальной размерности $n$ контактного многообразия $M^{2 n+1}$. Примером Л. м. в многообразии контактных элементов многообразия $N$ служит множество контак-





\end{document}